% CLAUDE REPORT::B3-ORACLE-0001::2026-11-13T22:50Z
% !TEX program = xelatex
% Compiles with XeLaTeX or pdfLaTeX (ASCII source).
\documentclass[11pt]{article}
\usepackage[margin=1in]{geometry}
\usepackage{amsmath,amssymb,amsthm}
\usepackage{xurl}
\newtheorem{theorem}{Theorem}
\newtheorem{proposition}[theorem]{Proposition}
\newtheorem{corollary}[theorem]{Corollary}
\newtheorem{lemma}[theorem]{Lemma}
\theoremstyle{definition}
\newtheorem{definition}[theorem]{Definition}
\theoremstyle{remark}
\newtheorem*{remark}{Remark}
\newcommand{\SAT}{\mathrm{SAT}}
\newcommand{\UNSAT}{\mathrm{UNSAT}}
\newcommand{\TQBF}{\mathrm{TQBF}}
\newcommand{\Wb}{\overline W}
\newcommand{\status}[1]{\textbf{[#1]}}

\title{CLAUDE B3-ORACLE-0001\\[2pt]
\large The weakest bridge B3 has no black-box or relativizing proof:\\
an oracle collapse theorem for $\SAT\in P \Rightarrow \Wb_R(N)\le\mathrm{poly}(N)$}
\author{CLAUDE (scheduled routine, worst-case derivation directive)}
\date{2026-11-13 14:50 America/Los\_Angeles}

\begin{document}
\maketitle

\noindent\textbf{Status.} P versus NP remains OPEN. This is an additive report under
\texttt{research/agents/claude/}. It changes no schedules, task count, role ownership, scientific
allocation, TMR/LSPV region, or root \texttt{main.tex}. It addresses the bridge B3 of the shared
worst-case directive
(\texttt{research/coordination/2026-11-02\_1204\_WORST\_CASE\_MINIMUM\_DERIVATION\_FIVE\_ROLE\_REFOCUS.md}).
Its main theorem is elementary; its value is to pin down, with proof, \emph{where} the
non-trivial content of the two-half program must sit.

\section{Setting and exact quantifiers}

Fix, independently of instances and sizes, a sound and complete refutation system $R$ for CNF
unsatisfiability: a polynomial-time decidable relation $R(F,\pi)$ (``$\pi$ is an $R$-derivation of
$\bot$ from $F$'') with a cost $\mathrm{cost}_R(\pi)\in\mathbb{N}_{\ge 1}$. As in the directive,
\[
W_R(F)=\min_{\pi:\,F\vdash_R\bot}\mathrm{cost}_R(\pi),\qquad
\Wb_R(N)=\max_{F\in\UNSAT,\ |\langle F\rangle|\le N} W_R(F)\quad(\max\emptyset=0).
\]
The inner minimum grants free knowledge of the best derivation (the user's
``minimal reduced logical state, minimal work'' over all valid reduction sequences); search cost is
never charged. Write
\[
X(R):\quad \exists C,k\ \forall N:\ \Wb_R(N)\le C(N+1)^k,
\qquad
\mathrm{LB}(R):=\neg X(R),
\qquad
\mathrm{B3}(R):\quad \SAT\in P\ \Rightarrow\ X(R).
\]
Call $R$ \emph{size-compatible} if there is a polynomial $q$ such that every $F$ with
$W_R(F)=w$ has an $R$-derivation $\pi$ with $|\pi|\le q(|\langle F\rangle|+w)$. (Unit-cost
DAG systems with polynomially bounded line size, e.g.\ Resolution, Extended Resolution with
paid definitions, and the project's 3-local mixed calculi, are size-compatible.)

The essential observation is that $X(R)$ \emph{mentions no algorithm and no oracle}: it is a
purely combinatorial statement about CNFs and $R$-derivations.

\section{Main results}

\begin{theorem}[Oracle collapse of B3]\label{thm:collapse}
\status{PROVED/EXACT}
For an oracle $O$ define the relativized bridge
$\mathrm{B3}^O(R):\ \SAT^O\in P^O\Rightarrow X(R)$, where $\SAT^O$ is satisfiability of
CNFs with $O$-query atoms (the standard relativization), and also the plain variant
$\mathrm{B3}_{\rm plain}^O(R):\ \SAT\in P^O\Rightarrow X(R)$. Then:
\begin{enumerate}
\item For $O=\TQBF$ (any PSPACE-complete set) both $\mathrm{B3}^O(R)$ and
$\mathrm{B3}_{\rm plain}^O(R)$ are \emph{equivalent to} $X(R)$.
\item For $O=\SAT$, $\mathrm{B3}_{\rm plain}^O(R)$ is equivalent to $X(R)$.
\item Hence if $\mathrm{LB}(R)$ holds, B3 is \emph{false} relative to $\TQBF$, and every
proof of $\mathrm{B3}(R)$ that relativizes is, in fact, a proof of $X(R)$.
\end{enumerate}
\end{theorem}

\begin{proof}
(1) $P^{\TQBF}=\mathrm{PSPACE}=NP^{\TQBF}$, so $\SAT^{\TQBF}\in NP^{\TQBF}=P^{\TQBF}$ and
$\SAT\in NP\subseteq P^{\TQBF}$. The hypotheses are true, so each implication holds iff its
conclusion $X(R)$ holds; $X(R)$ does not depend on $O$.
(2) $\SAT\in P^{\SAT}$ by one query, so again the implication is equivalent to $X(R)$.
(3) Immediate from (1): a relativizing argument proves $\mathrm{B3}^O(R)$ for every $O$,
in particular for $O=\TQBF$, i.e.\ it proves $X(R)$. If $\mathrm{LB}(R)$ is true, then
$\mathrm{B3}^{\TQBF}(R)$ is false.
\end{proof}

\begin{corollary}[What a relativizing B3 proof would cost]\label{cor:np}
\status{PROVED/EXACT}
If $R$ is size-compatible, a relativizing proof of $\mathrm{B3}(R)$ yields $NP=coNP$.
\end{corollary}
\begin{proof}
By Theorem~\ref{thm:collapse}(3), $X(R)$ holds, and size-compatibility gives every
unsatisfiable $F$ an $R$-refutation of length polynomial in $|\langle F\rangle|$, checkable in
polynomial time. So $\UNSAT\in NP$, i.e.\ $NP=coNP$ (Cook--Reckhow 1979).
\end{proof}

\begin{theorem}[Black-box indistinguishability]\label{thm:bb}
\status{PROVED/EXACT}
Call an argument for B3 \emph{black-box in the decider} if it uses a polynomial-time SAT
decider $A$ only through (a) its input--output function and (b) the existence of a polynomial
time bound for some oracle machine computing that function with oracle access to a set $O$.
Then any such argument proves $X(R)$ outright.
\end{theorem}
\begin{proof}
Every correct decider has input--output function $\chi_{\SAT}$, which exists independently of
whether $\SAT\in P$. The oracle machine ``query $O=\SAT$ on the input'' runs in linear time and
computes $\chi_{\SAT}$. So the premises (a) and (b) are satisfied unconditionally, and the
argument's conclusion $X(R)$ follows unconditionally.
\end{proof}

\begin{remark}[What this says about the program]
Theorems~\ref{thm:collapse}--\ref{thm:bb} locate the burden exactly. The lower-bound half
$\mathrm{LB}(R)$ is oracle-free, so a combinatorial proof of it (width, cut, potential,
hypergraph or ``minimal reduced state'' arguments with free optimal sequence selection) is
automatically relativization-proof but contributes \emph{no} non-relativizing content. All of the
non-relativizing content that any proof of $P\neq NP$ must contain (Baker--Gill--Solovay 1975)
must therefore be supplied by B3, and B3 must use the \emph{code} of the polynomial-time SAT
algorithm, not its behaviour. The only known mechanism of this kind is Cook's (1975)
tableau/PV translation of polynomial-time computations into Extended Frege proofs, which needs
a \emph{proof} of the algorithm's correctness, not merely its truth.
\end{remark}

\begin{corollary}[B3 is not logically weaker than the goal once LB is proved]\label{cor:equiv}
\status{PROVED/EXACT}
If $\mathrm{LB}(R)$ is a theorem then $\mathrm{B3}(R)\iff\SAT\notin P$.
In particular, for every $R$ for which a superpolynomial lower bound is already a
\textsc{known theorem} (Resolution: Haken 1985; Cutting Planes: Pudl\'ak 1997;
bounded-depth Frege: Ajtai 1988, Kraj\'{\i}\v{c}ek--Pudl\'ak--Woods and
Pitassi--Beame--Impagliazzo 1993/95; Polynomial Calculus: Razborov 1998;
regular Res($\oplus$): Efremenko--Garl\'{\i}k--Itsykson 2024), proving B3 is exactly proving
$P\neq NP$.
\end{corollary}
\begin{proof}
With $X(R)$ false, ``$\SAT\in P\Rightarrow X(R)$'' is equivalent to $\neg(\SAT\in P)$.
\end{proof}

\begin{proposition}[Known non-black-box sufficient route]\label{prop:pv}
\status{KNOWN THEOREM + CONDITIONAL}
If $R$ p-simulates Extended Frege (equivalently Extended Resolution) and
$\mathrm{cost}_R(\pi)\le|\pi|^{O(1)}$ for every derivation $\pi$, then
\[
\big(\SAT\in P\ \Rightarrow\ \mathrm{PV}\vdash P=NP\big)\ \Longrightarrow\ \mathrm{B3}(R).
\]
\end{proposition}
\begin{proof}[Proof sketch]
Cook--Kraj\'{\i}\v{c}ek (JSL 2007): $\mathrm{PV}\vdash P=NP$ iff PV defines a polynomial-time
algorithm returning, for each formula, a satisfying assignment or an EF proof of
unsatisfiability. Hence EF is p-bounded, and so is any system that p-simulates it, which gives
$X(R)$ because cost is polynomially bounded by length. The premise ``$\SAT\in P\Rightarrow\mathrm{PV}\vdash P=NP$'' is OPEN; by
Theorem~\ref{thm:collapse} it cannot be established by relativizing means for an $R$ satisfying
LB. This is the route already recorded in the corpus (LSPV-5/317, LSPV-1/361); it is restated
here only to mark it as the sole known non-black-box mechanism.
\end{proof}

\begin{corollary}[The price of a stronger $R$]\label{cor:er}
\status{PROVED/EXACT}
If $R$ p-simulates Extended Resolution (as the project's mixed calculi with fully paid
DEF/extension gates are recorded to do via Cook's ER shortcut controls), then $\mathrm{LB}(R)$
implies a superpolynomial Extended Resolution/EF lower bound, open since Cook--Reckhow 1979; by
Cook--Kraj\'{\i}\v{c}ek 2007 that would in turn imply that $P\neq NP$ is consistent with PV.
Conversely, choosing $R$ weak enough for $\mathrm{LB}(R)$ to be a known theorem makes
B3 equivalent to $P\neq NP$ (Corollary~\ref{cor:equiv}). There is no fixed $R$ for which
\emph{both} halves are presently within reach.
\end{corollary}
\begin{proof}
If $R$-cost is at most a polynomial in ER-size, a superpolynomial lower bound on $\Wb_R$ is a
superpolynomial lower bound on ER-size. The rest is the cited theorem and
Corollary~\ref{cor:equiv}.
\end{proof}

\section{Falsifiers stated before the work}
\begin{itemize}
\item[F1] A fixed size-compatible $R$ with $\mathrm{LB}(R)$ and a B3 proof that only uses the
input--output behaviour of the decider. \emph{Refuted} by Theorem~\ref{thm:bb}.
\item[F2] An oracle $O$ with $\SAT\in P^O$ under which B3 still holds for a known-hard $R$.
\emph{Refuted}: Theorem~\ref{thm:collapse}(1)--(2).
\item[F3] ``B3 is strictly weaker than $P\neq NP$ and therefore a cheaper target.'' True only as
long as LB is unproved; after LB it is equivalent (Corollary~\ref{cor:equiv}).
\item[F4] An $R$ whose LB is a known theorem and which also p-simulates EF. None known; finding
one would be a superpolynomial EF lower bound (Corollary~\ref{cor:er}).
\end{itemize}

\section{Post-quantum and barrier-transfer classification}
\begin{itemize}
\item PQ ledger item 5 (worst-case/average-case, externally extracting a resource from the
black-box behaviour of an arbitrary solver): \textbf{FAILED TRANSFER} for B3, with proof
(Theorem~\ref{thm:bb}). Any extraction must read the solver's code.
\item Proof complexity / logic (bounded arithmetic, Cook's PV translation): \textbf{directly
transferable theorem} (Proposition~\ref{prop:pv}), but conditional on a provability upgrade.
\item Relativization (Baker--Gill--Solovay) and algebrization (Aaronson--Wigderson 2009):
Theorem~\ref{thm:collapse} also covers algebrizing arguments, since with an algebraic extension
$\widetilde{\TQBF}$ we still have $\SAT\in P^{\widetilde{\TQBF}}$ while $X(R)$ is oracle-free.
\end{itemize}

\section{Corpus reviewed, duplication audit, changes}
\textbf{CORPUS REVIEWED.} Repository tree; \texttt{AGENTS.md}; the 2026-11-02 worst-case refocus
directive including LSPV-2/341 and LSPV-4 appendices; 2026-11-13 user two-gate handoff;
TMR replacement commit \texttt{e38c5f30} and TMR-1..5 directory listings; LSPV-1/357, 361, 367
and LSPV-5/317 lines on B3 found by repository search; prior CLAUDE world-ledger rows L01--L40.

\textbf{DUPLICATION / CONTRADICTION AUDIT.} The corpus repeatedly labels the missing B3 arrow
``nonrelativizing'' (LSPV-1/357, 367, 372, 589) and records for an artificial $R_{\rm tick}$ that
B3 is equivalent to $\SAT\notin P$ (LSPV-5/317). Repository search found no proof of the general
oracle-collapse statement or of the black-box indistinguishability theorem; Theorems
\ref{thm:collapse}--\ref{thm:bb} and Corollary~\ref{cor:equiv} supply them for every fixed $R$.
Proposition~\ref{prop:pv} duplicates the corpus's Cook-route observation and is credited, not
claimed. No contradiction found.

\textbf{WHAT CHANGED RELATIVE TO ROOT MAIN.} Nothing; \texttt{main.tex} not edited (this routine
owns no region).

\textbf{ADOPTED / REJECTED / RETIRED.} Adopted: directive definitions of $W_R,\Wb_R$, B1--B3.
Rejected as a route: any black-box or relativizing derivation of B3, and black-box solver
extraction (PQ item 5) for this arrow.

\textbf{STALE OR OVERSTATED CLAIMS FOUND.} Any phrasing that calls B3 the ``weakest'' bridge
should add that, after LB, B3 is equivalent to $\SAT\notin P$ and carries all of the
non-relativizing content.

\section{Literature checked this run}
Cook--Kraj\'{\i}\v{c}ek, ``Consequences of the provability of $NP\subseteq P/poly$'', JSL 72
(2007); Cook--Reckhow, JSL 44 (1979); Cook, STOC 1975 (PV); Baker--Gill--Solovay, SICOMP 1975;
Aaronson--Wigderson, TOCT 2009; Efremenko--Garl\'{\i}k--Itsykson, STOC 2024 (regular
Res($\oplus$)); Bhattacharya--Chattopadhyay, arXiv:2507.23008 (depth-restricted Res($\oplus$));
ECCC TR25-116/118 (2025 Res($\oplus$) work). Status: superpolynomial lower bounds for general
DAG-like Res($\oplus$) remain open, so even a calculus with GF(2) rows plus resolution is above
the current lower-bound frontier, consistent with Corollary~\ref{cor:er}.

\section{Next missing implication}
For the selected fixed $R$: a \emph{code-using} argument showing
$\SAT\in P\Rightarrow\mathrm{PV}\vdash$ (soundness of some polynomial-time SAT decider), or
another non-black-box route to $X(R)$; and, separately, $\mathrm{LB}(R)$. Both remain OPEN.
P versus NP remains OPEN.

\end{document}
% END_OF_VERIFIED_REPORT::CLAUDE::B3-ORACLE-0001
